% Module 1 section file. Included by ../module1_stellarator_equilibrium_geometry.tex.
\section{A reduced analytic equilibrium for intuition}
\label{sec:analytic_model}

\subsection{Circular nested surfaces}
Before using a three-dimensional equilibrium, it is useful to define a manufactured geometry with circular concentric surfaces:
\begin{align}
R(r,\theta)&=R_0+r\cos\theta,\\
Z(r,\theta)&=r\sin\theta,
\end{align}
where $0\leq r\leq a$. Let
\begin{equation}
s=\left(\frac{r}{a}\right)^2,
\qquad
r=a\sqrt{s}.
\end{equation}
This model is axisymmetric and therefore not a stellarator, but it provides an interpretable verification limit for coordinate mappings, Jacobians, volume, and interpolation.

The enclosed volume is approximately
\begin{equation}
V(r)=2\pi^2R_0r^2,
\end{equation}
and hence
\begin{equation}
V(s)=2\pi^2R_0a^2s,
\qquad
V'(s)=2\pi^2R_0a^2.
\end{equation}
These relations give exact or near-exact checks for a large-aspect-ratio circular implementation.

\subsection{Prescribed rotational transform}
A simple transform profile is
\begin{equation}
\iota(s)=\iota_0+\iota_1s,
\label{eq:iota_linear}
\end{equation}
where $\iota_0$ is the on-axis transform and $\iota_1$ is the total edge-minus-axis change. The shear measure becomes
\begin{equation}
\hat s_\iota(s)
=\frac{s\iota_1}{\iota_0+\iota_1s}.
\end{equation}
This profile is not derived from force balance; it is a manufactured geometry input suitable for checking downstream equations.

\subsection{Simple helical field-line tracing}
In straight-field-line coordinates,
\begin{equation}
\frac{d\theta}{d\zeta}=\iota(s).
\end{equation}
For constant $s$ and constant $\iota$,
\begin{equation}
\theta(\zeta)=\theta_0+\iota\zeta,
\label{eq:field_line_solution}
\end{equation}
where $\theta_0$ is the initial poloidal angle. This exact solution is a basic field-line tracer test.

\subsection{What the analytic model omits}
The circular manufactured model omits three-dimensional shaping, finite-beta equilibrium shifts, realistic current distributions, magnetic wells, coil ripple, islands, and stochasticity. It is a verification case, not a predictive stellarator equilibrium.

